\begin{examplebox}
	\textbf{\textcolor{CExample}{例 1（基础）}}

	\textbf{\textcolor{CExample}{题目：}} 已知长方体的长、宽、高分别为 $3,4,5$，求其外接球半径 $R$。

	\textbf{\textcolor{CExample}{解题步骤：}}
	\begin{description}[style=nextline, leftmargin=2.8em, labelsep=0.5em, itemsep=0.45em]
		\item[\textbf{第一步（套用公式）：}] 由于长方体满足条件，直接使用外接球半径公式。
		      \[
			      R = \frac{1}{2}\sqrt{a^2+b^2+c^2},\quad a=3,b=4,c=5.
		      \]
		      \par\textbf{理由：}
		      已知长方体的长、宽、高，可直接代入。

		\item[\textbf{第二步（计算平方和）：}] 计算平方和：
		      \[
			      3^2+4^2+5^2 = 9+16+25 = 50.
		      \]
		      \par\textbf{理由：}
		      体对角线长度的平方等于长、宽、高三个平方之和。

		\item[\textbf{第三步（求半径）：}]
		      \[
			      R = \frac{1}{2}\sqrt{50}
			      = \frac{1}{2}\cdot 5\sqrt{2}
			      = \frac{5\sqrt{2}}{2}.
		      \]
		      \par\textbf{理由：}
		      $\sqrt{50}=5\sqrt{2}$，外接球半径是体对角线的一半。
	\end{description}

	\textbf{\textcolor{CExample}{关键结论：}} \textbf{外接球半径 $R = \frac{5\sqrt{2}}{2}$。}

	\medskip
	\textbf{\textcolor{CExample}{例 2（稍变形）}}

	\textbf{\textcolor{CExample}{题目：}} 已知长方体的外接球表面积为 $36\pi$，长与宽之比为 $2:1$，高为 $3$，求长方体的长和宽。

	\textbf{\textcolor{CExample}{解题步骤：}}
	\begin{description}[style=nextline, leftmargin=2.8em, labelsep=0.5em, itemsep=0.45em]
		\item[\textbf{第一步（设未知量）：}] 设长为 $2x$，宽为 $x$，高为 $3$，其中 $x>0$。
		      \par\textbf{理由：}
		      长与宽之比为 $2:1$，所以可用同一个正数 $x$ 表示。

		\item[\textbf{第二步（利用表面积求半径）：}] 外接球表面积公式为
		      \[
			      S=4\pi R^2.
		      \]
		      由题意得
		      \[
			      4\pi R^2 = 36\pi \implies R^2 = 9 \implies R = 3.
		      \]
		      \par\textbf{理由：}
		      半径为正，所以取 $R=3$。

		\item[\textbf{第三步（代入半径公式列方程）：}] 由长方体外接球半径公式，
		      \[
			      \frac{1}{2}\sqrt{(2x)^2+x^2+3^2}=3.
		      \]
		      即
		      \[
			      \frac{1}{2}\sqrt{5x^2+9}=3
			      \implies
			      \sqrt{5x^2+9}=6.
		      \]
		      \par\textbf{理由：}
		      等式两边先同乘以 $2$，再准备平方。

		\item[\textbf{第四步（解方程求 $x$）：}] 两边平方，得
		      \[
			      5x^2+9=36.
		      \]
		      所以
		      \[
			      5x^2=27
			      \implies
			      x^2=\frac{27}{5}
			      \implies
			      x=\frac{3\sqrt{15}}{5}.
		      \]
		      因此
		      \[
			      \text{长}=2x=\frac{6\sqrt{15}}{5},
			      \qquad
			      \text{宽}=x=\frac{3\sqrt{15}}{5}.
		      \]
		      \par\textbf{理由：}
		      $x$ 表示长度，必须取正值。
	\end{description}

	\textbf{\textcolor{CExample}{关键结论：}} \textbf{长 $\dfrac{6\sqrt{15}}{5}$，宽 $\dfrac{3\sqrt{15}}{5}$。}

	\medskip
	\textbf{\textcolor{CExample}{例 3（综合应用）}}

	\textbf{\textcolor{CExample}{题目：}} 三棱锥 $A-BCD$ 中，$AB\perp$ 平面 $BCD$，$BC\perp CD$，$AB=2$，$BC=3$，$CD=4$。求三棱锥外接球半径 $R$。

	\textbf{\textcolor{CExample}{解题步骤：}}
	\begin{description}[style=nextline, leftmargin=2.8em, labelsep=0.5em, itemsep=0.45em]
		\item[\textbf{第一步（识别垂直结构）：}]
		      因为 $AB\perp$ 平面 $BCD$，所以 $AB$ 垂直于平面 $BCD$ 内的直线 $BC$ 和 $CD$。又因为 $BC\perp CD$，所以 $AB,BC,CD$ 对应的三个方向互相垂直。

		      \par\textbf{理由：}
		      这是把三棱锥补成长方体的关键：必须找到三个互相垂直的方向。

		\item[\textbf{第二步（坐标化补形）：}]
		      取 $B$ 为原点，令 $BC$ 沿 $x$ 轴，$CD$ 沿 $y$ 轴，$BA$ 沿 $z$ 轴。可设
		      \[
			      B(0,0,0),\quad C(3,0,0),\quad D(3,4,0),\quad A(0,0,2).
		      \]
		      这四点可以看作一个长方体的四个顶点，该长方体的三个方向长度分别为
		      \[
			      BC=3,\qquad CD=4,\qquad AB=2.
		      \]

		      \par\textbf{理由：}
		      虽然 $AB,BC,CD$ 不是从同一个顶点出发的三条棱，但它们对应三个互相垂直的方向，因此可放入长方体模型中处理。

		\item[\textbf{第三步（说明外接球相同）：}]
		      补形后的长方体外接球经过长方体所有顶点，因此也经过原三棱锥的四个顶点 $A,B,C,D$。又因为 $A,B,C,D$ 四点不共面，所以三棱锥的外接球唯一。

		      \par\textbf{理由：}
		      长方体外接球经过原三棱锥四个顶点，而四个不共面点唯一确定一个球，所以原三棱锥外接球与该长方体外接球相同。

		\item[\textbf{第四步（代入公式）：}]
		      因此原三棱锥外接球半径为
		      \[
			      R
			      =
			      \frac{1}{2}\sqrt{2^2+3^2+4^2}
			      =
			      \frac{1}{2}\sqrt{4+9+16}
			      =
			      \frac{\sqrt{29}}{2}.
		      \]
		      \par\textbf{理由：}
		      转化为长、宽、高分别为 $2,3,4$ 的长方体外接球问题。
	\end{description}

	\textbf{\textcolor{CExample}{关键结论：}} \textbf{外接球半径 $R = \dfrac{\sqrt{29}}{2}$。}
\end{examplebox}